\chapter{Introdução}
\label{cap:introducao}

\section{Contextualização}

A reabilitação motora de pessoas que perderam, parcial ou totalmente, a capacidade de
locomoção é um dos grandes desafios contemporâneos da área da saúde. Condições
neurológicas como o acidente vascular cerebral (AVC) e a lesão medular figuram entre as
principais causas de incapacidade física de longo prazo, comprometendo a marcha, a
independência e a qualidade de vida dos indivíduos acometidos \cite{rajashekar2024}. Há um
consenso crescente de que a recuperação funcional depende de treinamento motor
\emph{intensivo, repetitivo e engajador}, capaz de estimular os mecanismos de
reorganização do sistema nervoso \cite{loureiro2024}. É nesse contexto que tecnologias
assistivas como a robótica de reabilitação, a eletroestimulação funcional (FES) e a
realidade virtual (RV) vêm sendo progressivamente incorporadas à prática clínica e à
pesquisa acadêmica.

No Brasil, uma das iniciativas de referência nessa frente é o \textbf{Projeto EMA},
conduzido no Laboratório de Automação e Robótica (LARA) da Universidade de Brasília (UnB).
O projeto investiga o uso da eletroestimulação funcional integrada a dispositivos
robóticos --- como triciclos, esteiras, equipamentos de remo e andadores --- para a
reabilitação motora e cardiovascular de pessoas com lesão nos membros inferiores. Ao longo
de anos, sucessivas gerações de estudantes de graduação e pós-graduação
desenvolveram, nesse ecossistema, uma coleção de subsistemas de hardware e software: nós
de controle de estimuladores, drivers de sensores inerciais, interfaces de operação,
ambientes virtuais imersivos e ferramentas de aquisição e análise de dados. A
Figura~\ref{fig:ecossistema} apresenta uma visão geral desse ecossistema e situa os dois
subsistemas examinados neste trabalho.

\begin{figure}[H]
	\centering
	\caption{Visão geral do ecossistema do Projeto EMA e os dois estudos de caso deste trabalho.}
	\label{fig:ecossistema}
	\vspace{6pt}
	\begin{tikzpicture}[
		font=\small,
		bloco/.style={draw, rounded corners, align=center, inner sep=5pt, minimum height=0.95cm, minimum width=2.7cm},
		nucleo/.style={draw, rounded corners, align=center, inner sep=6pt, fill=black!8, minimum height=1cm, minimum width=3.4cm},
		caso/.style={draw, very thick, rounded corners, align=center, inner sep=5pt, minimum height=0.95cm, minimum width=2.7cm},
		]
		\node[nucleo] (ema) {Projeto EMA\\(FES + Robótica de\\Reabilitação --- LARA/UnB)};
		\node[caso, below=1.5cm of ema, xshift=-4.6cm] (trike) {Ciclismo com FES\\(\emph{trike})};
		\node[caso, below=1.5cm of ema, xshift=4.6cm] (vr) {Realidade Virtual\\(avatar imersivo)};
		\node[bloco, below=1.5cm of ema, xshift=-0.0cm, yshift=0.0cm, minimum width=3.1cm] (outros) {Remo, esteira,\\andador etc.};
		\node[bloco, below=1.2cm of trike] (hwt) {Estimulador, IMUs,\\Raspberry Pi, HMI};
		\node[bloco, below=1.2cm of vr] (hwv) {HMD, IMUs,\\\emph{game engine}};
		\foreach \n in {trike,outros,vr}{\draw[-{Latex[length=2mm]}] (ema.south) -- (\n.north);}
		\draw[-{Latex[length=2mm]}] (trike.south) -- (hwt.north);
		\draw[-{Latex[length=2mm]}] (vr.south) -- (hwv.north);
	\end{tikzpicture}

	\vspace{4pt}
	{\footnotesize Fonte: elaborado pelos autores.}
\end{figure}

Do ponto de vista da funcionalidade, muitos desses subsistemas cumpriram --- e alguns
ainda cumprem --- o seu propósito de pesquisa: chegaram a funcionar, geraram resultados e
sustentaram publicações. Entretanto, quando observados sob a ótica da
\emph{Engenharia de Software}, revelam um conjunto recorrente de fragilidades que
compromete a sua continuidade: arquiteturas monolíticas e fortemente acopladas, ausência
de documentação técnica e de guias de instalação, código sem padronização, repositórios de
versionamento desorganizados e uma lacuna sistemática entre aquilo que está
\emph{registrado} (em artigos, diagramas e trabalhos anteriores) e aquilo que está de fato
\emph{montado e em operação} no laboratório. O resultado é que cada novo integrante que
chega ao grupo enfrenta uma curva de reaproveitamento custosa, muitas vezes tendo de
reconstruir, por tentativa e erro, conhecimento que deveria estar documentado.

É precisamente esse o objeto deste Trabalho de Conclusão de Curso. Diferentemente de um
trabalho voltado a criar uma nova funcionalidade, a contribuição aqui proposta é de
natureza metodológica e diagnóstica: aplicar o instrumental da Engenharia de Software para
\emph{compreender, documentar, diagramar e diagnosticar} o estado atual dos sistemas do
Projeto EMA, apontando as suas limitações e propondo o caminho para que evoluam de
protótipos de pesquisa frágeis para uma base de software sustentável. Para dar concretude a
essa investigação, o trabalho organiza-se em torno de \textbf{dois estudos de caso}
complementares, que representam dois estágios distintos de maturidade dentro do mesmo
ecossistema.

\section{Os Dois Estudos de Caso}
\label{sec:estudos-de-caso}

O trabalho é desenvolvido em dupla e articula duas frentes que compartilham o mesmo
laboratório, o mesmo orientador e os mesmos problemas estruturais de Engenharia de
Software, mas que se encontram em estágios opostos do ciclo de vida:

\begin{itemize}
	\item \textbf{Estudo de caso A --- Realidade Virtual.} O sistema de RV imersivo
	      desenvolvido para capturar o movimento do paciente por sensores inerciais e
	      reproduzi-lo, em tempo real, em um avatar virtual \cite{balbino2020}. Trata-se de
	      um sistema cujo último desenvolvimento registrado no repositório data de 2019 e que
	      se encontra inativo desde então; o seu desafio central é a \emph{reativação}:
	      recuperar, a partir do que está registrado, um sistema que já funcionou,
	      identificando o que falta para colocá-lo novamente em operação. Por envolver menos
	      componentes de hardware, é o caso de mais rápida retomada.
	\item \textbf{Estudo de caso B --- Ciclismo com Eletroestimulação Funcional (FES
	      \emph{Cycling}).} O sistema de controle da \emph{trike} assistida por
	      eletroestimulação, no qual sensores inerciais capturam a posição dos pedais e
	      comandam os estimuladores que produzem o movimento de pedalada
	      \cite{brindeiro2017,baptista2022}. Diferentemente do caso anterior, é um sistema
	      \emph{ativo}, com equipe trabalhando sobre ele, porém marcado por documentação
	      ausente e por dificuldades de reprodução. Por envolver mais componentes de
	      hardware, é o caso de diagnóstico mais complexo.
\end{itemize}

A escolha por dois casos em estágios distintos é deliberada: permite contrastar os
problemas de Engenharia de Software que surgem quando um sistema é \emph{abandonado} com
aqueles que persistem mesmo quando ele está \emph{ativo}, revelando que a fragilidade
estrutural independe do estado de atividade do projeto. Ao longo do texto, os problemas
específicos de cada caso são tratados em capítulos próprios, enquanto os problemas comuns
--- notadamente os relativos ao versionamento e à documentação --- são consolidados em um
diagnóstico conjunto.

\section{Problema de Pesquisa}
\label{sec:problema}

Os sistemas do Projeto EMA foram concebidos, em sua maioria, como provas de conceito
orientadas a demonstrar viabilidade técnica e a sustentar investigações científicas
pontuais. Cumprido esse objetivo imediato, tais sistemas raramente receberam o tratamento
de engenharia necessário para que se tornassem artefatos reutilizáveis: faltam-lhes
documentação, padronização, modularidade e uma organização de versionamento que preserve o
conhecimento produzido. O maior gargalo do laboratório, portanto, não está na ausência de
funcionalidade, mas na \emph{forma como o software está estruturado, documentado e
versionado}.

Esse cenário produz uma consequência prática mensurável: a dificuldade de
\emph{reprodutibilidade}. Confrontado com um repositório do laboratório, um novo integrante
--- ainda que tecnicamente competente --- não encontra os elementos mínimos que a boa
prática de Engenharia de Software prescreve para que um sistema possa ser instalado,
executado e compreendido: um \emph{README} com passo a passo, a especificação do ambiente,
os diagramas de arquitetura coerentes com o que está de fato implementado e um histórico de
versões inteligível. Como se demonstra ao longo deste trabalho, a tentativa de
\emph{clonar e executar} os sistemas tal como se encontram esbarra em obstáculos que vão de
detalhes de configuração a lacunas estruturais, como código-fonte ausente e divergências
entre o hardware montado e o hardware documentado.

Diante desse quadro, o presente trabalho organiza-se em torno da seguinte questão de
pesquisa: \emph{de que forma o instrumental da Engenharia de Software pode ser aplicado ao
diagnóstico, à documentação e à reestruturação dos sistemas robóticos de reabilitação do
Projeto EMA, de modo a reduzir a lacuna entre o que está registrado e o que está montado e
a viabilizar que um novo integrante seja capaz de reproduzir, compreender e evoluir esses
sistemas?}

\section{Justificativa}
\label{sec:justificativa}

A relevância deste trabalho manifesta-se em três dimensões complementares.

Do ponto de vista da \textbf{Engenharia de Software} --- área de formação dos autores ---,
o laboratório oferece um cenário real, com legado e restrições concretos, no qual
princípios consagrados da disciplina podem ser aplicados e avaliados. A separação de
interesses, o baixo acoplamento, a documentação arquitetural e o versionamento disciplinado
são reconhecidos como determinantes da manutenibilidade e da longevidade de sistemas de
software \cite{sommerville2011}. Atuar sobre esses aspectos, com postura de \emph{arquiteto
de software} --- e não apenas de programador ---, é precisamente onde um trabalho da área
pode oferecer a maior contribuição a um laboratório historicamente centrado no hardware e
na pesquisa aplicada.

Do ponto de vista \textbf{científico}, o Projeto EMA é uma iniciativa de longa duração,
conduzida por sucessivas gerações de estudantes. A continuidade desse esforço depende
criticamente da capacidade de cada nova geração compreender e construir sobre o trabalho
das anteriores. Sem documentação e sem uma arquitetura inteligível, cada transição de
equipe implica retrabalho e perda de conhecimento --- um custo que a alta rotatividade
característica de um laboratório acadêmico apenas amplifica. Estruturar, documentar e
diagnosticar os sistemas é, assim, condição para a sustentabilidade das próprias linhas de
pesquisa.

Do ponto de vista \textbf{clínico}, os sistemas em questão são ferramentas com potencial
terapêutico documentado. A literatura demonstra que a realidade virtual e a
eletroestimulação funcional, sobretudo quando combinadas à robótica de reabilitação, podem
aumentar a motivação e a adesão ao tratamento e contribuir para melhores desfechos motores
\cite{loureiro2024,macchitella2023}. Recuperar e consolidar essas ferramentas em bases de
software confiáveis aproxima o laboratório de protocolos de reabilitação mais ricos, em
benefício direto dos pacientes.

\section{Objetivos}
\label{sec:objetivos}

\subsection{Objetivo Geral}

Aplicar o instrumental da Engenharia de Software ao diagnóstico, à documentação e à
proposta de reestruturação dos sistemas robóticos de reabilitação do Projeto EMA, tomando
como estudos de caso o sistema de realidade virtual e o sistema de ciclismo com
eletroestimulação funcional, de modo a reduzir a lacuna entre o que está registrado e o que
está montado e a viabilizar a sua reprodução, compreensão e evolução.

\subsection{Objetivos Específicos}

Para o alcance do objetivo geral, definem-se os seguintes objetivos específicos:

\begin{itemize}
	\item revisar a literatura sobre boas práticas de Engenharia de Software aplicáveis a
	      sistemas embarcados e robóticos, com ênfase em refatoração, documentação técnica,
	      versionamento e arquitetura baseada em componentes;
	\item levantar e documentar a arquitetura atual (\emph{as-is}) de cada um dos dois
	      sistemas, reconstruindo, por análise documental e de código, os seus componentes
	      e fluxos de dados;
	\item confrontar o que está registrado (artigos, diagramas e trabalhos anteriores) com
	      o que está de fato montado e em operação, mapeando as lacunas e as incoerências;
	\item diagnosticar as limitações de Engenharia de Software comuns aos dois sistemas,
	      em especial as relativas à organização dos repositórios, à documentação e à
	      reprodutibilidade;
	\item produzir os modelos e diagramas de arquitetura padronizados que representem os
	      sistemas de forma coerente;
	\item propor recomendações e um caminho de reestruturação --- de padronização de
	      repositórios e documentação a modularização e desacoplamento --- que preparem o
	      terreno para a etapa de integração a ser conduzida em trabalho posterior.
\end{itemize}

\section{Metodologia}
\label{sec:metodologia-intro}

Quanto à sua natureza, este trabalho caracteriza-se como uma \emph{pesquisa aplicada}, de
caráter \emph{exploratório} e abordagem predominantemente \emph{qualitativa}, voltada à
compreensão e à reestruturação de sistemas de software existentes em um contexto real de
pesquisa. Adota-se uma postura de \emph{consultoria de Engenharia de Software}: os autores
posicionam-se como engenheiros que, ao chegarem ao laboratório, recebem acesso aos
repositórios e à infraestrutura e se propõem a reproduzir os sistemas tal como se
encontram, registrando sistematicamente os obstáculos e as lacunas encontrados. O
detalhamento dessa abordagem é apresentado no Capítulo~\ref{cap:metodologia}.

\section{Organização do Trabalho}
\label{sec:organizacao}

Além desta introdução, o trabalho está organizado em mais sete capítulos. O
Capítulo~\ref{cap:fundamentacao} apresenta a \emph{fundamentação teórica}, começando pelas
boas práticas de Engenharia de Software para sistemas embarcados --- refatoração,
documentação técnica, versionamento e arquitetura baseada em componentes (ROS) --- e
delimitando, em seguida, os fundamentos de domínio (reabilitação, realidade virtual,
eletroestimulação e sensoriamento inercial). O Capítulo~\ref{cap:metodologia} descreve a
\emph{metodologia da consultoria} adotada para o diagnóstico. Os
Capítulos~\ref{cap:caso-vr} e \ref{cap:caso-trike} apresentam, respectivamente, os
\emph{estudos de caso} da realidade virtual e do ciclismo com eletroestimulação, com a
documentação da arquitetura atual, as tentativas de reprodução e as limitações específicas
de cada um. O Capítulo~\ref{cap:diagnostico} consolida o \emph{diagnóstico dos problemas
comuns}, com destaque para a organização do versionamento e a documentação. O
Capítulo~\ref{cap:proposta} reúne as \emph{recomendações e o caminho de reestruturação}
propostos. Por fim, o Capítulo~\ref{cap:consideracoes} traz as \emph{considerações finais},
retomando as contribuições, as limitações e os trabalhos futuros.
